\begin{table}[H]
\caption{Operacionalización de las variables}
\label{tab:operacionalizacion}
\setstretch{1}
\renewcommand{\arraystretch}{1.3}
\small
\begin{tabularx}{\textwidth}{|L{3.2cm}|L{3.4cm}|Y|}
\hline
\rowcolor{gray!25}
\textbf{Variable} & \textbf{Dimensiones} & \textbf{Indicadores} \\
\hline
1. VI. Modelo crítico Critic-EmpRAG1.0
 & 1.1. Mecanismo de recuperación (RAG)
 & 1.1.1. Estrategia de recuperación (densa, híbrida o adaptativa por incertidumbre) \newline
   1.1.2. Condición de activación de la recuperación (siempre activa vs.\ condicional) \\
\cline{2-3}
 & 1.2. Mecanismo de autorreflexión
 & 1.2.1. Tipo de autocrítica (verbal, \textit{tokens} aprendidos o basada en reglas) \newline
   1.2.2. Número máximo de ciclos de refinamiento iterativo \newline
   1.2.3. Fuente de verificación de la crítica (evidencia externa recuperada vs.\ conocimiento paramétrico) \\
\hline
2. VD. Productividad y acceso al conocimiento
 & 2.1. Fundamentación y confiabilidad
 & 2.1.1. Exactitud de respuesta (\textit{accuracy}/F1 frente a la respuesta esperada) \newline
   2.1.2. Tasa de alucinación (afirmaciones no sustentadas por la evidencia recuperada, sobre el total de afirmaciones verificables) \newline
   2.1.3. Proporción de afirmaciones con fuente o cita verificable \newline
   2.1.4. Puntaje de fidelidad (\textit{faithfulness}) respecto a los documentos recuperados \newline
   2.1.5. Tasa de abstención correcta ante consultas sin respaldo en el corpus \\
\cline{2-3}
 & 2.2. Productividad del proceso de consulta
 & 2.2.1. Tiempo medio de resolución por consulta operativa \newline
   2.2.2. Tasa de consultas resueltas sin escalamiento a personal experimentado \newline
   2.2.3. Interrupciones semanales a personal experimentado por consultas de conocimiento \newline
   2.2.4. Esfuerzo percibido en la consulta (escala Likert de 1 a 5) \\
\cline{2-3}
 & 2.3. Viabilidad bajo restricciones
 & 2.3.1. Latencia promedio de respuesta \newline
   2.3.2. Consumo de recursos de cómputo (memoria/GPU) \newline
   2.3.3. Proporción de consultas resueltas sin transmisión de datos a servicios de terceros \newline
   2.3.4. Superficies de exposición con control implementado, sobre el total identificado \\
\hline
\end{tabularx}

\vspace{0.2cm}
\noindent{\small \textit{Nota.} Los indicadores de la dimensión 2.2 se registran en los contextos $A_0$ y $A_1$ con el mismo instrumento, conforme al diseño preexperimental; los de las dimensiones 2.1 y 2.3 se obtienen sobre el banco de evaluación en las configuraciones de la Tabla \ref{tab:configuraciones}.}
\end{table}
